\documentclass[preview]{standalone}
\usepackage{amsmath}
\usepackage{booktabs}
\usepackage{amssymb}
\usepackage{setspace}
\usepackage{bm}
\let\oldcdot\cdot
\usepackage{breqn}
\let\cdot\oldcdot
\begin{document}
Let $\phi(x)$ be the level-set function and
 \begin{align*}\raisebox{\depth}{$\chi$}(x) &= \begin{cases} 1 & \mbox{if}\quad \phi(x) < 0 \\ k_{air}/k_0 & \mbox{if}\quad \phi(x)  \ge 0 \end{cases},\end{align*}
where the weak-material contrast is the physically motivated ratio of the thermal conductivity of air to that of aluminium, $k_{air}/k_0 = 0.026/200 \approx 1.3 \times 10^{-4}$ (numerically $1.3 \times 10^{-4}$, the value taken by the characteristic function).
The design domain is the normalized unit square — the cross-section of the spreader plate of a liquid-cooled power module scaled to unit size for benchmark comparison.  A power chip occupying the central forty percent of the bottom edge injects a uniform heat flux of one thousand watts per square metre; the remaining bottom edge is adiabatic and the bottom edge is a fixed boundary (the chip position is prescribed).  The TOP edge — opposite the chip — is the wall of a cooling pipe (a cold plate, i.e. a liquid-cooled cold rail): it is held at the prescribed reference temperature of zero kelvin, the zero of the temperature scale, an isothermal heat sink of negligible thermal resistance.  The left and right edges exchange heat with ambient air at $u_{amb} = 298$ K through an effective convective heat transfer coefficient of $h = 25$ W/(m$^{2}$ K).  The thermal conductivity of the aluminium spreader is TEMPERATURE-DEPENDENT: $k(u) = k_0 \left( 1 + \beta \left( u - u_{amb}\right)\right)$ with $k_0 = 200$ W/(m K) and the linear temperature coefficient $\beta = 5 \times 10^{-4}$ K$^{-1}$, a realistic order of magnitude for aluminium between 270 K and 340 K (the conductivity varies by $\pm 2\%$ across the working temperature span).  The $k(u)$ dependence makes the state problem NONLINEAR and is the physical mechanism studied here.  The optimisation minimises the chip temperature integral at fixed material cost.
The Lagrangian of the cold-plate problem is
\begin{dmath*}
\mathcal L(\left[ u\right],\left[ \delta u\right]) = J(u) -  \int_{\Gamma_q} q  \delta u \,ds_x +  \int_{D} k(u) \raisebox{\depth}{$\chi$} \left(  \nabla u  \cdot  \nabla \delta u  \right) \,dx + \int_{\Gamma_c} h \left( u - u_{amb}\right) \delta u \,ds_x ,
\end{dmath*}
with
\begin{align*}
J(u) &=  \int_{D} \ell \raisebox{\depth}{$\chi$} \,dx + \int_{\Gamma_q} q u \,ds_x ,\\
k(u) &= k_0 \left( 1 + \beta \left( u - u_{amb}\right)\right),
 \end{align*}
where the chip temperature integral — the thermal compliance that measures the thermal resistance of the spreader — is weighted by the injected flux $q = 1000$ W/m$^{2}$.  The Lagrange multiplier $\ell$ weights the material cost; it is a balance parameter of the optimisation, calibrated by the pipeline (a geometric scan locating the knee of the volume-fraction curve), not part of the problem statement.
The state $\left[ u\right] \in \mathbb P_{2}(D)$ is the solution of the NON-LINEAR problem
\begin{dmath*}
\int_{D} k(u) \raisebox{\depth}{$\chi$} \left(  \nabla u  \cdot  \nabla \delta u  \right) \,dx + \int_{\Gamma_c} h u \delta u \,ds_x = \int_{\Gamma_q} q  \delta u \,ds_x + \int_{\Gamma_c} h u_{amb} \delta u \,ds_x \qquad \forall\; {\delta u\in \mathbb P_{2}(D)} ,
\end{dmath*}
with the Dirichlet value $\mathtt{\text{0.}}$ on the cold plate.
The problem is solved by Newton's method: the state $u$ is updated as $u \leftarrow u - \hat u$, where the Newton update $\left[ \hat u\right]$ at the state $\left[ u\right]$ is the solution of
\begin{dmath*}
\int_{D} k(u) \raisebox{\depth}{$\chi$} \left(  \nabla \hat u  \cdot  \nabla \delta u  \right) \,dx + \int_{D} k_0 \beta \hat u \raisebox{\depth}{$\chi$} \left(  \nabla u  \cdot  \nabla \delta u  \right) \,dx + \int_{\Gamma_c} h \hat u \delta u \,ds_x = \int_{D} k(u) \raisebox{\depth}{$\chi$} \left(  \nabla u  \cdot  \nabla \delta u  \right) \,dx + \int_{\Gamma_c} h \left( u - u_{amb}\right) \delta u \,ds_x - \int_{\Gamma_q} q  \delta u \,ds_x \qquad \forall\; {\delta u\in \mathbb P_{2}(D)} ,
\end{dmath*}
where the conductivity term is linearised as $\partial_{u} \left[ k(u) \nabla u \cdot \nabla \delta u\right] \hat u = k(u) \nabla \hat u \cdot \nabla \delta u + k_0 \beta \hat u \nabla u \cdot \nabla \delta u$.
The adjoint state $\left[ \tilde u\right]$ to the direct state $\left[ u\right]$ is the solution of
\begin{dmath*}
\int_{D} k(u) \raisebox{\depth}{$\chi$} \left(  \nabla \delta u  \cdot  \nabla \tilde u  \right) \,dx + \int_{D} k_0 \beta \delta u \raisebox{\depth}{$\chi$} \left(  \nabla u  \cdot  \nabla \tilde u  \right) \,dx + \int_{\Gamma_c} h \delta u \tilde u \,ds_x = - \int_{\Gamma_q} q \delta u \,ds_x \qquad \forall\; {\delta u\in \mathbb P_{2}(D)} ,
\end{dmath*}
where the adjoint operator is the linearisation of the state operator at $u$ (the non-symmetric conductivity linearisation appears in its transposed form), and the adjoint load $- \int_{\Gamma_q} q \delta u \,ds_x$ comes from the partial derivative $\partial_{u} J(\delta u) = \int_{\Gamma_q} q \delta u \,ds_x$ of the chip temperature integral in the objective.
In order to compute the shape derivative, we consider a shape perturbation $\mathbf x = T_t(\mathbf X) = \mathbf X + t \mathbf  V(\mathbf X)$ with a suitable velocity field $\mathbf V$, vanishing on the fixed chip footprint and the cold plate.  The perturbed Lagrangian $\mathcal G$ using the Lagrangian state $\left[ u\right]^{L}$ is given by
\begin{dmath*}
\mathcal G(t,\left[ u\right]^{L},\left[ \delta u\right]^{L}) = \int_{D} \ell \raisebox{\depth}{$\chi$}  \operatorname{det}\mathbf F( t )  \,dx + \int_{\Gamma_q} q u^{L} \,ds_x -  \int_{\Gamma_q} q  \delta u^{L} \,ds_x +  \int_{D} k(u) \raisebox{\depth}{$\chi$}  \operatorname{det}\mathbf F( t )  \left( \left(  \nabla u^{L} \cdot  \mathbf F^{-1}(t) \right) \cdot \left(  \nabla \delta u^{L} \cdot  \mathbf F^{-1}(t) \right) \right) \,dx + \int_{\Gamma_c} h \left( u^{L} - u_{amb}\right) \delta u^{L} \,J_{s}(t) \,ds_x ,
\end{dmath*}
where $\mathbf F(t) = \mathbf I + t \nabla \mathbf V$ is the deformation gradient and $J_{s}(t)$ the boundary Jacobian of the moving cooling edges $\Gamma_c$; the fixed chip footprint $\Gamma_q$ carries no Jacobian because $\mathbf V$ vanishes there, and the temperature-dependent conductivity $k(u)$ depends on the STATE but not on the shape, so it is treated as a coefficient under the shape perturbation.
For the direct state $\left[ u\right]$ and the adjoint state $\left[ \tilde u\right]$, the volume expression of the shape derivative, including the curvature term on the moving cooling edges, is given by
\begin{dmath*}
DJ(\Omega)(\mathbf V) = \int_{D} \left(\operatorname{div} {\mathbf V} \left(  \nabla u  \cdot  \nabla \tilde u  \right) - \left(  \nabla \tilde u  \cdot \left(  \nabla u  \cdot  \nabla {\mathbf V}  \right) \right) - \left(  \nabla u  \cdot \left(  \nabla \tilde u  \cdot  \nabla {\mathbf V}  \right) \right)\right) k(u) \raisebox{\depth}{$\chi$} \,dx + \int_{D} \ell \raisebox{\depth}{$\chi$} \operatorname{div} {\mathbf V} \,dx + \int_{\Gamma_c} \kappa h \left( u - u_{amb}\right) \tilde u \left(\mathbf V \cdot \mathbf n\right) \,ds_x ,
\end{dmath*}
where $\kappa$ is the curvature of $\Gamma_c$ and $\mathbf n$ its outward unit normal.  The chip footprint and the cold plate are fixed boundaries; the cooling edges $\Gamma_c$ are design boundaries — their shape is optimised.
The initial design is an I-beam without holes: a full-width top flange (thickness 0.2) lying along the cooling wall, a bottom flange (length 0.5, thickness 0.2) under the chip footprint, and a web of width 0.2 running the full height between them.  The temperature field uses P2 finite elements with the cooling-wall temperature applied directly to the wall segment.
\end{document}
